\elegant_provide_module:n{math}


\tl_if_exist:NF \g__elegant_math_tl
  {
    \tl_new:N \g__elegant_math_tl
    \tl_gset:Nn \g__elegant_math_tl { cm }
  }

\clist_new:N \g__elegant_math_clist

\elegant_msg_new:nn { not support math }
  {
    The~math~option~`#1'~is~not~support.\\
    Please~select~one~of~the~following~math~options.\\
    \clist_use:Nn \g__elegant_math_clist {,~}
  }


\cs_new_protected_nopar:Npn \elegant_declare_math:nn #1#2
  {
    \hook_new:n { elegant-module / math / #1 / before }
    \hook_new:n { elegant-module / math / #1 / after }
    \tl_gset:cn { g__elegant_math_#1_tl } { #2 }
    \clist_gput_right:Nn \g__elegant_math_clist { #1 }
  }


\cs_new_protected_nopar:Npn \elegant_apply_selected_math:
  {
    \tl_if_exist:cTF { g__elegant_math_ \g__elegant_math_tl _tl }
      {
        \hook_use:n { elegant-module / math / \g__elegant_math_tl / before }
        \tl_use:c { g__elegant_math_ \g__elegant_math_tl _tl }
        \hook_use:n { elegant-module / math / \g__elegant_math_tl / after }
      }
      {
        \elegant_msg_error:nx { not support math } { \g__elegant_math_tl }
      }
  }

\cs_new_protected_nopar:Npn \__elegant_math_load_unicode_math:
  {
    \RequirePackage { unicode-math }
    \keys_set:nn { unicode-math }
      {
        math-style = ISO,
        bold-style = ISO,
      }
  }


\elegant_declare_math:nn { cm } { \relax }

\elegant_declare_math:nn { newtx }
  {
    \let \oldencodingdefault \encodingdefault
    \let \oldrmdefault \rmdefault
    \let \oldsfdefault \sfdefault
    \let \oldttdefault \ttdefault
    \def \encodingdefault { T1 }
    \renewcommand { \rmdefault } { ntxtlf }
    \renewcommand { \sfdefault } { qhv }
    \renewcommand { \ttdefault } { ntxtt }
    \RequirePackage { newtxmath }
    \let \encodingdefault \oldencodingdefault
    \let \rmdefault \oldrmdefault
    \let \sfdefault \oldsfdefault
    \let \ttdefault \oldttdefault
    \let \Bbbk \relax
    \RequirePackage { esint }

    %%% use yhmath pkg, uncomment following code
    % \let \oldwidering \widering
    % \let \widering \undefined
    % \RequirePackage { yhmath }
    % \let \widering \oldwidering

    %%% use esvect pkg, uncomment following code
    % \RequirePackage { esvect }

    \DeclareSymbolFont { CMlargesymbols } { OMX } { cmex } { m } { n }
    \let \sumop \relax
    \let \prodop \relax
    \DeclareMathSymbol { \sumop } { \mathop } { CMlargesymbols } { "50 }
    \DeclareMathSymbol { \prodop } { \mathop } { CMlargesymbols } { "51 }
  }

\elegant_declare_math:nn { mtpro2 }
  {
    \let \Bbbk \relax
    \RequirePackage [ lite ] { mtpro2 }
  }

\elegant_declare_math:nn { stixtwo }
  {
    \__elegant_math_load_unicode_math:
    \setmathfont { STIXTwoMath-Regular.otf }
    \setmathfont { STIXTwoMath-Regular.otf }
      [ range = { scr, bfscr }, StylisticSet = 01 ]
  }

\elegant_declare_math:nn { xits }
  {
    \__elegant_math_load_unicode_math:
    \setmathfont { XITSMath-Regular.otf }
      [ BoldFont = XITSMath-Bold.otf ]
    \setmathfont { XITSMath-Regular.otf }
      [ range = { cal, bfcal }, StylisticSet = 01 ]
  }

\elegant_declare_math:nn { tg }
  {
    \__elegant_math_load_unicode_math:
    \setmathfont { texgyretermes-math.otf }
    \setmathfont { XITSMath-Regular.otf }
      [ range = { cal, bfcal }, StylisticSet = 01 ]
  }

\elegant_apply_selected_math:

\RequirePackage { fixdif }


\endinput
